\documentclass[11pt,a4paper]{article}
\usepackage[utf8]{inputenc}
\usepackage[vietnamese]{babel}
\usepackage{amsmath,amssymb}
\usepackage{geometry}
\geometry{a4paper, total={180mm,260mm}, left=15mm, right=15mm, top=18mm, bottom=18mm}

\usepackage{tikz}
\usetikzlibrary{shapes.geometric, arrows.meta, positioning, calc}
\tikzset{
  arr/.style={-{Stealth[scale=0.8]}},
  every node/.append style={align=center}
}
\usepackage{tcolorbox}
\tcbuselibrary{skins,breakable}
\usepackage{fontawesome}
\usepackage{fancyhdr}
\usepackage{enumitem}

\pagestyle{fancy}
\fancyhf{}
\renewcommand{\headrulewidth}{0.5pt}
\renewcommand{\footrulewidth}{0.5pt}
\fancyhead[L]{\small\bfseries\color{violet!85!black}\faBook\ SỔ TAY KIẾN THỨC TOÁN 2}
\fancyhead[R]{\small\bfseries\color{violet!85!black}Lớp Toán Thầy Lĩnh -- loptoanthaylinh.vn}
\fancyfoot[C]{\small\color{gray}-- Sổ tay kiến thức: Bài 31. Thực hành xem đồng hồ, xem lịch --}

\setlength{\parindent}{0pt}
\setlength{\parskip}{0pt}

\begin{document}

% --- BANNER BÀI HỌC ---
\begin{tcolorbox}[
  enhanced,
  colback=violet!8!white,
  colframe=violet!70!black,
  arc=3mm, boxrule=1.2pt,
  left=3mm, right=3mm, top=2mm, bottom=2mm,
  halign=center
]
  {\LARGE\bfseries\color{violet!90!black}\faBookmark\ BÀI 31: THỰC HÀNH XEM ĐỒNG HỒ, XEM LỊCH}\\[1mm]
  {\itshape\footnotesize\color{black!75} Kĩ năng thực tế: Đọc đồng hồ điện tử, tính thời lượng hoạt động và tra cứu lịch trên 1 trang A4}
\end{tcolorbox}

\vspace{-2mm}

\noindent
\begin{minipage}[t]{0.49\textwidth}\vspace{0pt}
\begin{tcolorbox}[
  enhanced,
  colback=violet!4!white,
  colframe=violet!75!black,
  arc=2mm, boxrule=0.8pt,
  left=2mm, right=2mm, top=1.8mm, bottom=1.8mm,
  title={\footnotesize\bfseries\color{white}\faKey\ 1. Đọc đồng hồ điện tử}
]
{\footnotesize
\begin{itemize}\setlength{\itemsep}{1pt}\setlength{\leftskip}{-2.5mm}
  \item Đồng hồ điện tử hiển thị dạng: \textbf{Giờ : Phút}.
  \item Hai chữ số trước dấu hai chấm chỉ \textbf{giờ}, hai chữ số sau chỉ \textbf{phút}.
  \item \textbf{Ví dụ:}
  \begin{itemize}\setlength{\itemsep}{0.5pt}\setlength{\leftskip}{-2.5mm}
    \item $08:00 \implies 8$ giờ đúng sáng.
    \item $14:30 \implies 14$ giờ $30$ phút ($2$ giờ rưỡi chiều).
    \item $20:00 \implies 20$ giờ ($8$ giờ tối).
  \end{itemize}
\end{itemize}
}
\end{tcolorbox}
\end{minipage}
\hfill
\begin{minipage}[t]{0.49\textwidth}\vspace{0pt}
\begin{tcolorbox}[
  enhanced,
  colback=teal!4!white,
  colframe=teal!75!black,
  arc=2mm, boxrule=0.8pt,
  left=2mm, right=2mm, top=1.8mm, bottom=1.8mm,
  title={\footnotesize\bfseries\color{white}\faSliders\ 2. Tính khoảng thời gian}
]
{\footnotesize
\begin{itemize}\setlength{\itemsep}{1pt}\setlength{\leftskip}{-2.5mm}
  \item \textbf{Thời gian trôi qua:}\\
  Muốn biết một hoạt động diễn ra trong bao lâu, ta lấy \textbf{giờ kết thúc} trừ đi \textbf{giờ bắt đầu}.
  $$\mathbf{\text{Thời lượng} = \text{Giờ kết thúc} - \text{Giờ bắt đầu}}$$
  \item \textbf{Ví dụ:} Nam bắt đầu học lúc $8$ giờ, học xong lúc $10$ giờ. Nam học trong:\\
  $10 - 8 = \mathbf{2\text{ giờ}}$.
\end{itemize}
}
\end{tcolorbox}
\end{minipage}

\vspace{-1.5mm}
\begin{tcolorbox}[
  enhanced,
  colback=orange!4!white,
  colframe=orange!80!black,
  arc=2mm, boxrule=0.8pt,
  left=2mm, right=2mm, top=1.5mm, bottom=1.5mm,
  title={\footnotesize\bfseries\color{white}\faSitemap\ 3. Sơ đồ liên kết: Đồng hồ kim và Đồng hồ điện tử}
]
\centering
\begin{tikzpicture}[x=1.5cm, y=0.8cm]
  \node[rectangle, rounded corners=2pt, draw=violet!80!black, fill=violet!15, text width=4.5cm, align=center, font=\scriptsize\bfseries] at (0,0) {
    \textbf{ĐỒNG HỒ KIM}\\Kim ngắn chỉ giữa $3$ và $4$\\Kim dài chỉ số $6$
  };
  \draw[<->, very thick, orange!80!black] (1.8,0) -- node[above, font=\tiny\bfseries] {Cùng chỉ} (3.2,0);
  \node[rectangle, rounded corners=2pt, draw=teal!80!black, fill=teal!15, text width=4.5cm, align=center, font=\scriptsize\bfseries] at (5,0) {
    \textbf{ĐỒNG HỒ ĐIỆN TỬ}\\Sáng: $03:30$ ($3$ giờ $30$ phút)\\Chiều: $15:30$ ($3$ giờ rưỡi chiều)
  };
\end{tikzpicture}
\end{tcolorbox}

\vspace{-1.5mm}
\noindent
\begin{minipage}[t]{0.49\textwidth}\vspace{0pt}
\begin{tcolorbox}[
  enhanced,
  colback=green!4!white,
  colframe=green!60!black,
  arc=2mm, boxrule=0.8pt,
  left=2mm, right=2mm, top=1.8mm, bottom=1.8mm,
  title={\footnotesize\bfseries\color{white}\faCheckCircle\ 4. Ví dụ thực hành chuẩn mực}
]
{\footnotesize
\textbf{Bài toán thực tế:}\\
Một buổi xem phim bắt đầu lúc $19$ giờ và kết thúc lúc $21$ giờ. Hỏi buổi xem phim kéo dài bao nhiêu giờ?\\
\textit{Bài giải:}\\
Buổi xem phim kéo dài số giờ là:
$$21 - 19 = 2\text{ (giờ)}$$
Đáp số: $2$ giờ.
}
\end{tcolorbox}
\end{minipage}
\hfill
\begin{minipage}[t]{0.49\textwidth}\vspace{0pt}
\begin{tcolorbox}[
  enhanced,
  colback=red!4!white,
  colframe=red!75!black,
  arc=2mm, boxrule=0.8pt,
  left=2mm, right=2mm, top=1.8mm, bottom=1.8mm,
  title={\footnotesize\bfseries\color{white}\faExclamationTriangle\ 5. Mẹo tránh bẫy trắc nghiệm}
]
{\footnotesize
\begin{itemize}\setlength{\itemsep}{1pt}\setlength{\leftskip}{-2.5mm}
  \item \textbf{Cẩn thận khi trừ giờ sáng và chiều:}\\
  Nếu bắt đầu lúc $11$ giờ trưa, kết thúc lúc $2$ giờ chiều: Phải đổi $2$ giờ chiều thành $14$ giờ, rồi lấy $14 - 11 = 3$ giờ. Không lấy $11 - 2$!
  \item Luôn quan sát kĩ kim ngắn xem đã qua mốc số mấy.
\end{itemize}
}
\end{tcolorbox}
\end{minipage}


\end{document}
